\chapter{Optimization I: fewer launches and kernel fusion}\label{ch:fusion}

The port's first rule produces many kernels: each loop nest of the source becomes its own kernel, in the source order
(\cref{ch:port}). This keeps the arithmetic recognizably the source's, and it is the right way to get correct. It is
not the fastest way to run. This chapter shows how the number of launches is reduced afterwards without changing a
bit: by merging kernels that run the same statements, by fusing producers and consumers, by letting independent
kernels overlap, and by replaying recorded launch sequences.

\section{What a launch costs}

A kernel launch costs time in three places:
\begin{itemize}
  \item on the host, where the OpenACC runtime checks that the arrays are present, packs the arguments and submits
    the launch;
  \item on the device, where the kernel's blocks are distributed to the SMs, and where the last blocks of a kernel run
    while most SMs already idle (the \emph{tail});
  \item between kernels, where a kernel cannot start until the previous one in its queue has finished.
\end{itemize}
The profiler measures these as the gaps between kernels (\cref{ch:prof}). For a kernel that runs 100~$\mu$s they do
not matter; for one that runs 5~$\mu$s they can double its cost. The model of \cref{ch:perf} shows that the full
case's d02 is mostly bandwidth-bound, while the dev case's d02 and the boundary strips of both are launch-bound.

\section{Fusion in bitwise terms}

Two kernels $K_1$ then $K_2$ over the same iteration space can be fused into one kernel that, at each point $p$,
executes $K_1$'s statements at $p$ and then $K_2$'s statements at $p$. The fused kernel computes the same values if
every value each statement reads is the value it read before. Let $W_1, R_1$ be the sets of (array, index) pairs that
$K_1$ writes and reads, and likewise $W_2, R_2$. The fusion is legal when:
\begin{enumerate}
  \item $K_2$ reads $K_1$'s outputs only at the same point: if $(a, q) \in R_2$ and $(a, q) \in W_1$, then the read
    is made at $p = q$. (Otherwise $K_2$ at $p$ might run before $K_1$ at $q$ has written.)
  \item $K_2$ does not overwrite what $K_1$ reads at another point: if $(a, q) \in W_2$ and $(a, q) \in R_1$, then
    $p = q$.
  \item Values both write at the same point are written in the original order, which the fused body keeps.
\end{enumerate}
Under these conditions every statement reads the same operands and performs the same operations: the fusion is in
class R1 (\cref{tab:bitclass}). A \emph{pointwise} producer and consumer always satisfy them. A \emph{stencil}
consumer, which reads the producer's output at neighbouring points, does not; it can be fused only in two ways:
\begin{itemize}
  \item each thread \emph{recomputes} the producer's values at the neighbouring points it needs, with the same
    operations on the same operands. The result is the same bits, at the cost of repeated work;
  \item a tile of the producer's output is computed into shared memory with its halo, synchronized, and then read
    (\cref{ch:tiling}).
\end{itemize}

When the two kernels have different iteration ranges, as the $u$ and $v$ updates do on the staggered grid, the fused
kernel runs over the union of the ranges and guards each statement group with its own range: the same technique as
Template C.

\begin{algorithm}[ht]
\caption{Fusion legality of two consecutive kernels (analyzer A1, on the read and write sets of the kernel bodies).}
\label{alg:fusion}
\begin{algorithmic}[1]
\Require kernels $K_1$, $K_2$ consecutive in one queue, no host statement between them
\For{each array $a$ written by $K_1$ and read by $K_2$}
  \If{$K_2$ reads $a$ at an offset $\ne 0$} \Return stencil: fuse only with recomputation or a tile \EndIf
\EndFor
\For{each array $a$ written by $K_2$ and read by $K_1$}
  \If{$K_1$ reads $a$ at an offset $\ne 0$} \Return illegal \EndIf
\EndFor
\If{the iteration ranges differ} \State fuse over the union, guard each body with its own range \EndIf
\State \Return legal (R1); estimated saving: one launch, plus the write and re-read of each intermediate array that
  is not needed afterwards
\end{algorithmic}
\end{algorithm}

Fusion saves more than a launch when an intermediate array is used only to pass values from $K_1$ to $K_2$: the fused
kernel can keep them in registers and never write them. That changes no arithmetic either: the value in a register is
the value that would have been stored.

\section{Merging boundary strips}

Template G turns each strip of a boundary routine into its own kernel (\cref{ch:port}). The strips of one call are
disjoint and execute the same statement on different points. Merging the four strips of a call into one kernel, with
an index map from the merged iteration space to the strip and point, changes nothing but the schedule: class R0
(optimization O2, phase S6-04). Strips of \emph{different phases} (the $x$ phase and the $y$ phase of
\code{set_physical_bc3d}) stay separate, because the second reads corners the first wrote.

\section{Asynchronous queues}

Kernels that touch disjoint data could run at the same time, but by default each waits for the previous one. With
\code{async(q)}, kernels go into queues, and kernels in different queues may overlap. Examples in WRF:
\begin{itemize}
  \item the boundary updates of $\theta$, $\mu$ and $\mu_t$ after \code{advance_mu_t}, which touch three different
    fields;
  \item the per-species work of scalar transport, where each species is independent until the next RK stage;
  \item the boundary-condition calls of different fields at the end of a stage.
\end{itemize}
The schedule changes, the statements do not: class R0. The risk is a missing \code{wait} before a kernel that needs
the results, which is a race and would show as a difference in T-AB. The dependency map comes from the read and write
sets that analyzer A1 already extracts (phase S6-05).

\section{CUDA Graphs}

A CUDA Graph records a sequence of launches once and replays it with a single call, which removes most of the host
cost per kernel. The acoustic sub-step is a natural candidate: the same kernels in the same order, seven times per
d02 step. Two difficulties make it an experiment (phase S6-08):
\begin{itemize}
  \item OpenACC does not expose graphs. They are recorded by CUDA stream capture on the CUDA stream of an OpenACC
    queue (\code{acc_get_cuda_stream}), from CUDA Fortran;
  \item a graph records the kernels' arguments, but some scalars change between sub-steps (the iteration number, the
    sub-step's time step in the first RK stage). They must be passed through device memory, or the recorded parameters
    updated before each replay.
\end{itemize}
A graph changes nothing in what the kernels compute: class R0.

\section{Worked example: the acoustic sub-step}

On d02 the acoustic sub-step calls, in order (\code{dyn_em/solve_em.F}):
\code{advance_uv}; the boundary updates of $u$ and $v$; \code{advance_mu_t}; the boundary updates of $\theta$,
$\mu$ and $\mu_t$; \code{advance_w}; \code{sumflux}; the boundary updates of $\phi$ and $w$; \code{calc_p_rho}.
From the port's kernel table, that is
\begin{itemize}
  \item 8 interior kernels: two in \code{advance_uv}, three in \code{advance_mu_t}, one each in \code{advance_w},
    \code{sumflux} and \code{calc_p_rho};
  \item 28 boundary-strip kernels: four strips for each of the six \code{spec_bdyupdate} calls ($u$, $v$, $\theta$,
    $\mu$, $\mu_t$, $w$) and four for \code{spec_bdyupdate_ph}.
\end{itemize}
That is 36 launches per sub-step and 252 per d02 step from the acoustic loop alone, most of them for strips five points
wide. The planned reductions, each with its class:
\begin{center}
\small
\begin{tabular}{p{0.52\textwidth}lr}
\toprule
Change & Class & Launches per sub-step \\
\midrule
as ported & --- & 36 \\
four strips of each boundary call merged (S6-04) & R0 & 15 \\
the three boundary calls after \code{advance_mu_t} merged into one multi-field strip kernel & R0 & 13 \\
the $u$ and $v$ kernels of \code{advance_uv} fused over the union of their ranges (S6-07) & R1 & 12 \\
\bottomrule
\end{tabular}
\end{center}
Further fusions need a legality note per chain, which phase S6-07 writes before any code. For example,
\code{advance_mu_t}'s first and third kernels both work column by column, and the third reads the vertical mass flux
the first wrote only in the same column; whether the second, pointwise kernel between them allows the fusion is
decided by \cref{alg:fusion}. The measured times before and after are \tbr.

\begin{portnote}
Fusion is the most tempting optimization to get wrong, because the fused code looks almost identical. The port's
safeguards are the same as for porting: a written legality note, \code{arith_guard} and \code{check_verbatim} on the
statements, T-AB of every route the fused kernel serves, and T-TRACE over the windows. T-AB must still be able to
isolate each routine, so a fused kernel needs a fallback: one simple design runs the unfused kernels whenever any of
the fused routes is switched off.
\end{portnote}
